% Heston vs Black–Scholes: Fitting the Volatility Smile — research-note PDF.
%
% PROSE AND CODE ARE VERBATIM from the website article:
%   site/src/content/articles/heston-vs-black-scholes.tsx
%   site/src/content/articles/heston-code.ts
% Metadata from lib/articles.ts; project-card fallbacks from lib/topics.ts.
%
% NO CHARTS. The article's header records it as deliberately unexecuted: no
% data module, no computed numbers, prose and code only.
\documentclass{pyportfolios-note}

\noteshorttitle{Heston vs Black–Scholes}

\begin{document}
\thispagestyle{notecover}

\notemasthead{Quantitative Insights}{Quant Insights}{Q3 2026}{}

\notetitle{Heston vs Black–Scholes:\\Fitting the Volatility Smile}

\notedeck{One model prices every strike with the same number. The other doesn't
have to.}

\textbf{The claim we test:} Black-Scholes assumes a single, constant volatility
— so it predicts the \textit{same} implied vol at every strike. The market
flatly disagrees: implied vol curves into a \textbf{smile/skew}. The Heston
model, by letting volatility itself be random, can bend to fit that curve. This
piece pits the two against a real SPX-style implied-vol surface and measures who
fits, by how much, and at what cost.

\notecard
  {Quant Foundations \& Derivatives}
  {QuantLib · SciPy · NumPy}
  {SPX / SPY options}
  {2019 – 2024 + recent option chain}
  {Pricing options consistently across strikes — the vol desk's answer to the
   smile BS can't fit}

\begin{notecode}
import numpy as np
import matplotlib.pyplot as plt

np.random.seed(42)
plt.rcParams["figure.dpi"] = 110
\end{notecode}

% ==================================================== 1 The disagreement =====
\noteneedroom{170bp}
\notesection{1.\notenumsep The disagreement}

\begin{multicols}{2}
Black-Scholes takes one volatility $\sigma$ and returns one price per strike.
Invert real option prices back into implied vols and — if BS were right — you'd
get a \textbf{flat line}: the same $\sigma$ at every strike. Instead, equity
index options show a pronounced \textbf{skew}: deep out-of-the-money puts trade
at much higher implied vol than at-the-money or upside calls.

Why? Two BS assumptions break:

\begin{notebullets}
\item \textbf{Volatility isn't constant} — it clusters, spikes, and
      mean-reverts.
\item \textbf{Returns aren't Gaussian} — crashes are fatter and more frequent
      than a normal distribution allows, and they cluster in falling markets.
\end{notebullets}

The skew is the market pricing in exactly those two facts. Any single-vol model
is structurally unable to reproduce it.
\end{multicols}

% ====================================================== 2 The contenders =====
\noteneedroom{200bp}
\notesection{2.\notenumsep The contenders}

\begin{notetable}{@{}p{130bp}p{160bp}p{195bp}@{}}
\thcell{} & \thcell{Black-Scholes} & \thcell{Heston} \\
\theadrule
Volatility & \textbf{Constant} $\sigma$ &
  \textbf{Stochastic} — its own mean-reverting process \\ \rowrule
Free parameters & 1 & 5 \\ \rowrule
Returns distribution & Log-normal (Gaussian log-returns) &
  Fat-tailed, skewed \\ \rowrule
Can fit the smile? & No — flat by construction & Yes \\ \rowrule
Closed form? & Yes (simple) &
  Semi-closed (characteristic function) \\
\end{notetable}

\textbf{Black-Scholes} assumes $dS = \mu S\,dt + \sigma S\,dW$ — the GBM from
Tutorial 1.

\textbf{Heston} adds a second equation for variance $v_t$:

\[ dS_t = \mu S_t\,dt + \sqrt{v_t}\,S_t\,dW_t^S, \qquad
   dv_t = \kappa(\theta - v_t)\,dt + \xi\sqrt{v_t}\,dW_t^v \]

with $\text{corr}(dW^S, dW^v) = \rho$. The five parameters: $\kappa$
(mean-reversion speed), $\theta$ (long-run variance), $\xi$ (vol-of-vol),
$\rho$ (spot-vol correlation, the skew driver), and $v_0$ (initial variance).

% ============================================================= 3 The test ====
\noteneedroom{170bp}
\notesection{3.\notenumsep The test}

We build a representative SPX implied-vol skew (ATM $\approx$ 20\%, steep put
wing — the shape index options actually trade) and give \textbf{both} models the
same job: match it.

\begin{multicols}{2}
\begin{notebullets}
\item \textbf{Black-Scholes} gets one degree of freedom — its best move is a
      single flat line, so we set it to the ATM vol.
\item \textbf{Heston} gets calibrated: we search its five parameters to minimise
      the gap to the market smile.
\end{notebullets}

\textit{Note: we use a representative parametric skew so the notebook runs
without a live option-chain subscription; the calibration machinery is identical
for a real SPX/SPY chain.}
\end{multicols}

\begin{notecode}
# --- representative SPX 3M implied-vol skew (moneyness K/S -> implied vol) ---
S0, T, r, q = 100.0, 0.25, 0.03, 0.0
moneyness = np.array([0.80, 0.85, 0.90, 0.95, 1.00, 1.05, 1.10, 1.15, 1.20])
strikes   = moneyness * S0

# market smile: high vol in the put wing, gentle rise in far calls (typical equity skew)
mkt_iv = np.array([0.32, 0.285, 0.255, 0.228, 0.205, 0.192, 0.188, 0.191, 0.198])

fig, ax = plt.subplots(figsize=(9, 5))
ax.plot(moneyness, mkt_iv * 100, "o-", color="black", label="Market implied vol (SPX-style skew)")
ax.axhline(mkt_iv[4] * 100, color="steelblue", ls="--", label="Black-Scholes (single vol = ATM)")
ax.set_xlabel("Moneyness  K / S"); ax.set_ylabel("Implied volatility (%)")
ax.set_title("The disagreement: market skew vs the flat line BS must draw")
ax.legend(); plt.tight_layout(); plt.show()
\end{notecode}

Black-Scholes is a horizontal line. The market is a curve. That gap — biggest in
the crash-protection put wing — is the mispricing a single-vol model bakes into
every out-of-the-money option.

\notesubsection{3.1\notenumsep Calibrate Heston (QuantLib pricing + bounded optimisation)}

A single-maturity smile under-determines five parameters — unconstrained
optimisers happily wander into absurd values ($\kappa$ in the thousands) that fit
the curve for the wrong reasons. So we calibrate with \textbf{economically
sensible bounds}: mean-reversion $\kappa \in [0.5, 8]$, vol-of-vol
$\xi \leq 2$, correlation $\rho \in [-0.95, -0.05]$, variances $\leq 0.25$.
QuantLib prices; SciPy searches.

\begin{notecode}
import QuantLib as ql
from scipy.optimize import least_squares

today = ql.Date(15, 6, 2024)
ql.Settings.instance().evaluationDate = today
day_count = ql.Actual365Fixed()
calendar = ql.NullCalendar()

spot_h = ql.QuoteHandle(ql.SimpleQuote(S0))
r_ts = ql.YieldTermStructureHandle(ql.FlatForward(today, r, day_count))
q_ts = ql.YieldTermStructureHandle(ql.FlatForward(today, q, day_count))
maturity = ql.Period(3, ql.Months)
expiry = today + maturity

def heston_ivs(params):
    """Price every strike under Heston, return Black-Scholes implied vols."""
    v0, kappa, theta, xi, rho = params
    proc = ql.HestonProcess(r_ts, q_ts, spot_h, v0, kappa, theta, xi, rho)
    eng = ql.AnalyticHestonEngine(ql.HestonModel(proc))
    bsm = ql.BlackScholesMertonProcess(spot_h, q_ts, r_ts,
            ql.BlackVolTermStructureHandle(
                ql.BlackConstantVol(today, calendar, 0.20, day_count)))
    ivs = []
    for K in strikes:
        opt = ql.VanillaOption(ql.PlainVanillaPayoff(ql.Option.Call, K),
                               ql.EuropeanExercise(expiry))
        opt.setPricingEngine(eng)
        try:
            ivs.append(opt.impliedVolatility(opt.NPV(), bsm))
        except Exception:
            ivs.append(np.nan)
    return np.array(ivs)

def residuals(p):
    return np.nan_to_num(heston_ivs(p) - mkt_iv, nan=0.5)

x0     = [0.042, 2.0, 0.045, 0.8, -0.7]                          # v0, kappa, theta, xi, rho
bounds = ([0.005, 0.5, 0.005, 0.05, -0.95], [0.25, 8.0, 0.25, 2.0, -0.05])
res = least_squares(residuals, x0, bounds=bounds, xtol=1e-12, ftol=1e-12)

v0, kappa, theta, xi, rho = res.x
print("Calibrated Heston parameters (bounded):")
print(f"  v0    (initial var)    = {v0:.4f}  -> vol {np.sqrt(v0):.1%}")
print(f"  theta (long-run var)   = {theta:.4f}  -> vol {np.sqrt(theta):.1%}")
print(f"  kappa (mean-reversion) = {kappa:.2f}")
print(f"  xi    (vol-of-vol)     = {xi:.2f}")
print(f"  rho   (spot-vol corr)  = {rho:.2f}   <- negative = downside skew")
\end{notecode}

\notesubsection{3.2\notenumsep Recover Heston's implied-vol smile}

Price each strike under the calibrated model and invert to implied vol — same
axis as the market curve.

\begin{notecode}
heston_iv = heston_ivs(res.x)
bs_iv = np.full_like(mkt_iv, mkt_iv[4])   # flat at ATM

print("moneyness  market   heston    BS(flat)")
for m, mk, he, bs in zip(moneyness, mkt_iv, heston_iv, bs_iv):
    print(f"  {m:.2f}     {mk:6.1%}  {he:6.1%}   {bs:6.1%}")
\end{notecode}

% ======================================================= 4 The scoreboard ====
\noteneedroom{150bp}
\notesection{4.\notenumsep The scoreboard}

\begin{notecode}
fig, (ax1, ax2) = plt.subplots(1, 2, figsize=(13, 5))

# left: the three curves
ax1.plot(moneyness, mkt_iv * 100, "o-", color="black", label="Market", lw=2)
ax1.plot(moneyness, heston_iv * 100, "s--", color="crimson", label="Heston (calibrated)")
ax1.plot(moneyness, bs_iv * 100, "--", color="steelblue", label="Black-Scholes (flat)")
ax1.set_xlabel("Moneyness  K / S"); ax1.set_ylabel("Implied vol (%)")
ax1.set_title("Fit to the smile"); ax1.legend()

# right: absolute pricing error in vol points
bs_err = np.abs(bs_iv - mkt_iv) * 100
he_err = np.abs(heston_iv - mkt_iv) * 100
x = np.arange(len(moneyness)); w = 0.38
ax2.bar(x - w/2, bs_err, w, color="steelblue", label="Black-Scholes")
ax2.bar(x + w/2, he_err, w, color="crimson", label="Heston")
ax2.set_xticks(x); ax2.set_xticklabels([f"{m:.2f}" for m in moneyness], rotation=45)
ax2.set_xlabel("Moneyness  K / S"); ax2.set_ylabel("|error| (vol points)")
ax2.set_title("Absolute implied-vol error by strike"); ax2.legend()
plt.tight_layout(); plt.show()

print(f"Mean abs error — Black-Scholes: {bs_err.mean():.2f} vol pts")
print(f"Mean abs error — Heston:        {he_err.mean():.2f} vol pts")
print(f"Worst strike (put wing) — BS: {bs_err[0]:.2f} pts  vs  Heston: {he_err[0]:.2f} pts")
\end{notecode}

The picture is decisive where it matters most — the \textbf{put wing}.
Black-Scholes underprices crash protection by a wide margin (its flat line sits
far below the market's elevated put vols), while Heston tracks the curve because
its negative $\rho$ generates exactly that downside skew. Across the surface
Heston's mean error is a fraction of Black-Scholes'.

% ============================================================= 5 Verdict =====
\noteneedroom{300bp}
\notesection{5.\notenumsep Verdict}

\textbf{Heston wins the fit — decisively — and it isn't close in the wings.}
With five parameters versus one, that's expected; the real question is whether
the extra complexity earns its keep. The answer depends on the job:

\begin{multicols}{2}
\begin{notebullets}
\item \textbf{Pricing a single vanilla at one strike?} Black-Scholes with that
      strike's \textit{own} implied vol is simpler and exact. The smile is a
      lookup table; you don't need a model to read one row.
\item \textbf{Pricing a book across many strikes consistently, or anything
      exotic (barriers, cliquets, forward-starts) whose value depends on the
      \textit{dynamics} of vol?} Heston is close to essential — a flat vol
      would misprice the smile-sensitive payoff.
\end{notebullets}

\textbf{What each model really is:}

\begin{notebullets}
\item Black-Scholes isn't wrong so much as \textit{incomplete} — it's the
      quoting convention (implied vol) more than a dynamic model. Its flat smile
      is a feature of using one number, not a bug to be fixed.
\item Heston buys smile-consistency and fat tails at the cost of five parameters
      that must be \textbf{recalibrated} as the surface moves, and a fit that
      can still miss very short maturities (where jumps, not stochastic vol,
      drive the steep skew).
\end{notebullets}
\end{multicols}

\textbf{The honest trade-off:} more parameters always fit better in-sample.
Heston earns trust because its parameters are \textit{economically
interpretable} ($\rho$ is the skew, $\xi$ the smile convexity, $\theta$ the vol
term structure) — it's fitting the smile \textit{for the right reasons}, not
just curve-fitting.

\textbf{Where to go next:} for the very short-dated skew Heston struggles with,
add jumps (\textbf{Bates}, Heston + Merton jumps); to fit the \textit{entire}
observed surface exactly rather than approximately, \textbf{Dupire local
volatility}; and for the constant-vol baseline both are measured against,
revisit the \textbf{Black-Scholes} tutorial.

\end{document}
